\section{A balanced moving-slope squarefree-energy theorem}
\label{sec:tpc33-squarefree-energy}

The normal form separates a support question from a sign question.  We now
show that the squarefree support has the expected positive density in a
genuine growing-slope range.  The proof is an elementary squarefree sieve
tailored to the B\'ezout coordinates.  The origin must be balanced: an
arbitrary transformation \((x,y)\mapsto(x+ak,y+sk)\) can make the heights
of both forms arbitrarily large while \(a,s\), and the length of the
parameter interval remain fixed.
For general background on elementary squarefree sieves, see
\citet{IwaniecKowalski2004}.

\subsection{Balanced coordinates and the Euler product}

Assume \eqref{eq:tpc33-affine-hypotheses}.  If \(a\ge2\), choose the unique
integer \(x\) such that
\begin{equation}
 sx\equiv1\pmod a,\qquad -\frac a2<x\le\frac a2,
 \label{eq:tpc33-balanced-x}
\end{equation}
and put \(y=(sx-1)/a\).  If \(a=1\), take \(x=0\), \(y=-1\).  Then
\begin{equation}
 sx-ay=1,\qquad |x|\le\frac a2,\qquad
 |y|\le\frac s2+1.
 \label{eq:tpc33-balanced-bounds}
\end{equation}
Use the forms \(D_t,U_t\) from
\eqref{eq:tpc33-canonical-affine-forms} and put
\begin{equation}
 K=\max(a,s).
 \label{eq:tpc33-K}
\end{equation}
For \(T\ge4(|h|+1)\) and \(T<t\le2T\), both forms are positive and
\begin{equation}
 \max(D_t,U_t)\ll_h KT.
 \label{eq:tpc33-affine-height}
\end{equation}

Define
\begin{equation}
 E_{a,s,h}(T)
 :=\sum_{\substack{T<t\le2T\\(t,h)=1}}
 \mu^2(D_t)\mu^2(U_t)
 \label{eq:tpc33-squarefree-energy}
\end{equation}
and
\begin{equation}
 \boxed{
 \mathfrak S_{\rm sf}(a,s;h)
 :=\prod_{p\mid h}\left(1-\frac1p\right)
 \prod_{\substack{p\nmid h\\p\mid as}}
       \left(1-\frac1{p^2}\right)
 \prod_{p\nmid ash}\left(1-\frac2{p^2}\right).}
 \label{eq:tpc33-squarefree-singular-series}
\end{equation}

\begin{theorem}[Uniform squarefree energy]
\label{thm:tpc33-uniform-squarefree-energy}
For fixed \(h\ne0\), uniformly over all positive \(a,s,T\) satisfying
\eqref{eq:tpc33-affine-hypotheses} and \(T\ge4(|h|+1)\),
\begin{equation}
 \boxed{
 E_{a,s,h}(T)
 =\mathfrak S_{\rm sf}(a,s;h)T
 +O_h\!\left(
   \frac{T}{\log(2T)}+\sqrt{KT}
  \right).}
 \label{eq:tpc33-squarefree-asymptotic}
\end{equation}
Furthermore,
\begin{equation}
 \mathfrak S_{\rm sf}(a,s;h)\ge c_h>0
 \label{eq:tpc33-squarefree-positive}
\end{equation}
uniformly in \(a,s\).  Consequently, for every fixed \(\eta>0\),
\begin{equation}
 K\le T^{1-\eta}
 \quad\Longrightarrow\quad
 E_{a,s,h}(T)
 =\mathfrak S_{\rm sf}(a,s;h)T+o_{h,\eta}(T)
 \label{eq:tpc33-squarefree-growing-slope-range}
\end{equation}
uniformly in that range.
\end{theorem}

\subsection{Local root counts}

\begin{lemma}[The three local factors]
\label{lem:tpc33-local-squarefree-factors}
The local density of the conditions in
\eqref{eq:tpc33-squarefree-energy} is:
\begin{enumerate}[label=\textup{(\roman*)}]
 \item \(1-1/p\) if \(p\mid h\);
 \item \(1-1/p^2\) if \(p\nmid h\) and \(p\mid as\);
 \item \(1-2/p^2\) if \(p\nmid ash\).
\end{enumerate}
\end{lemma}

\begin{proof}
If \(p\mid h\), then \(p\nmid as\).  The primitive condition excludes
\(t\equiv0\pmod p\), and
\[
 D_t\equiv st\pmod p,
 \qquad
 U_t\equiv at\pmod p.
\]
Thus neither form is divisible by \(p\) on the retained classes, giving
\(1-1/p\).

Suppose \(p\nmid h\) and \(p\mid s\).  Then \(p\nmid a\), and
\(sx-ay=1\) makes \(y\not\equiv0\pmod p\).  Hence \(D_t\equiv hy\not\equiv0
\pmod p\), while \(U_t\) has exactly one root modulo \(p^2\).  The case
\(p\mid a\) is symmetric.  This proves (ii).

Finally suppose \(p\nmid ash\).  Each affine form has exactly one root
modulo \(p^2\).  The roots are distinct: a common root would imply
\(p^2\mid sU_t-aD_t=h\), which is impossible.  Thus precisely two classes
modulo \(p^2\) are removed.  The argument includes \(p=2\).
\end{proof}

\subsection{Small primes and the large-square tail}

\begin{proof}[Proof of \cref{thm:tpc33-uniform-squarefree-energy}]
For \(z\ge2\), put
\begin{equation}
 M_z=\rad|h|
 \prod_{\substack{p\le z\\p\nmid h}}p^2.
 \label{eq:tpc33-small-prime-modulus}
\end{equation}
Let \(N_z(T)\) count \(T<t\le2T\) for which \((t,h)=1\) and neither
form is divisible by \(p^2\) for any \(p\le z\), \(p\nmid h\).  The
conditions are periodic modulo \(M_z\).  By the Chinese remainder theorem
and \cref{lem:tpc33-local-squarefree-factors},
\begin{equation}
 N_z(T)=\mathfrak S_z(a,s;h)T+O(M_z),
 \label{eq:tpc33-truncated-sieve-count}
\end{equation}
where \(\mathfrak S_z\) is the product in
\eqref{eq:tpc33-squarefree-singular-series} truncated at primes at most
\(z\), together with all fixed factors \(p\mid h\).  Absolute convergence
gives, uniformly in \(a,s\),
\begin{equation}
 |\mathfrak S_z(a,s;h)-\mathfrak S_{\rm sf}(a,s;h)|
 \ll_h\sum_{p>z}\frac1{p^2}
 \ll_h\frac1z.
 \label{eq:tpc33-singular-series-tail}
\end{equation}

It remains to remove integers for which a prime \(p>z\) has
\(p^2\mid D_tU_t\).  If \(p\mid s\), the congruence \(p^2\mid D_t\)
has no solution; otherwise it occupies at most one class modulo \(p^2\).
The analogous statement holds for \(U_t\).  By
\eqref{eq:tpc33-affine-height}, a relevant prime satisfies
\(p\ll_h\sqrt{KT}\).  Therefore
\begin{align}
 |N_z(T)-E_{a,s,h}(T)|
 &\ll_h
 \sum_{z<p\ll_h\sqrt{KT}}
 \left(\frac{T}{p^2}+1\right)
 \notag\\
 &\ll_h\frac{T}{z}+\sqrt{KT}.
 \label{eq:tpc33-large-square-tail}
\end{align}

For all sufficiently large \(T\), Chebyshev's estimate
\(\sum_{p\le z}\log p\ll z\) permits the choice
\(z=c\log(2T)\), with a sufficiently small absolute \(c>0\), for which
\(M_z\ll_h T^{1/2}\); the remaining bounded range is absorbed by the
implied constant.  Combining
\eqref{eq:tpc33-truncated-sieve-count}--
\eqref{eq:tpc33-large-square-tail}, and absorbing the harmless
\(T^{1/2}\) term, proves \eqref{eq:tpc33-squarefree-asymptotic}.

For positivity, replacing a factor \(1-2/p^2\) by
\(1-1/p^2\) can only increase the product.  Hence
\begin{equation}
 \mathfrak S_{\rm sf}(a,s;h)
 \ge
 \prod_{p\mid h}\left(1-\frac1p\right)
 \prod_{p\nmid h}\left(1-\frac2{p^2}\right)
 =:c_h>0.
 \label{eq:tpc33-squarefree-lower-product}
\end{equation}
Finally, if \(K\le T^{1-\eta}\), the second error in
\eqref{eq:tpc33-squarefree-asymptotic} is at most
\(T^{1-\eta/2}\), proving
\eqref{eq:tpc33-squarefree-growing-slope-range}.
\end{proof}

\begin{remark}[What the theorem does not give]
\label{rem:tpc33-squarefree-not-signed}
The summand in \eqref{eq:tpc33-squarefree-energy} is nonnegative.  The
theorem proves that squarefree values persist uniformly in the stated
moving-slope window; it gives no cancellation in
\(\sum\mu(D_t)\mu(U_t)\).  In the physical high-\(\beta\) columns the slope
can be much larger than the affine-fiber length, so even the asymptotic
range \(K\le T^{1-\eta}\) need not apply.
\end{remark}

\begin{remark}[Why the dyadic origin matters]
For a general interval \(T_0<t\le T_0+Y\), the same tail argument naturally
contains
\(\sqrt{K(|T_0|+Y)}\), not merely \(\sqrt{KY}\).  The balanced dyadic
formulation is therefore a genuine hypothesis rather than a cosmetic
normalization.
\end{remark}
